% =============================================================
% CHAPITRE 1 — CADRE DU PROJET
% =============================================================
\chapter{Cadre du projet}

\begin{chapintro}
Ce premier chapitre situe notre projet dans son contexte professionnel
et technique. Nous y présentons le domaine de la gouvernance des agents
IA, analysons l'existant pour en identifier les lacunes, décrivons la
solution AegisWeb, puis exposons la méthodologie et le planning retenus.
\end{chapintro}

\section{Présentation du contexte et de l'organisme d'accueil}

\subsection{Contexte du projet}

\textbf{AegisWeb} est une plateforme de \textbf{gouvernance et de confiance
pour agents IA agissant sur le web}. Elle fournit une couche de contrôle
entre les agents autonomes et les systèmes métier (portails SaaS,
formulaires, facturation, etc.). Son principe fondateur est
\textit{«~Permission before autonomy~»}~: la permission est le goulot
d'étranglement, pas l'intelligence.

Le premier cas d'usage ciblé est le \textbf{procurement SaaS}~: permettre
à des agents IA de gérer en toute sécurité les factures, les
renouvellements et les changements de plan soumis à approbation humaine,
sans exposer de credentials bruts.

\begin{table}[H]
\centering
\caption{Fiche d'identité du projet AegisWeb}
\label{tab:fiche_projet}
\begin{tabularx}{\textwidth}{@{}lX@{}}
\toprule
\textbf{Critère}         & \textbf{Description} \\
\midrule
Dénomination             & AegisWeb (nom interne historique~: AgentPass) \\
\rowcolor{rowgray}
Secteur d'activité       & Cybersécurité applicative, gouvernance IA \\
Nature                   & Plateforme MVP locale de gouvernance d'agents web \\
\rowcolor{rowgray}
Produit phare            & Couche permission, credential, approbation, runtime
                           navigateur et audit pour agents IA \\
Marché cible             & Entreprises utilisant des agents IA pour
                           l'automatisation web contrôlée \\
\rowcolor{rowgray}
Cas d'usage MVP          & Procurement SaaS (factures, renouvellements,
                           downgrades avec approbation) \\
Langues de l'interface   & Anglais (dashboard), documentation bilingue \\
\bottomrule
\end{tabularx}
\end{table}

\subsection{Cadre du projet de fin d'études}

Notre projet de fin d'études consiste en la \textbf{conception et le
développement complets de la plateforme AegisWeb}, couvrant~:
\begin{itemize}[leftmargin=1.8cm, itemsep=3pt]
  \item le \textbf{backend} (API REST NestJS, base de données, files
    BullMQ, endpoints internes worker) ;
  \item le \textbf{worker} (exécution de workflows, Playwright, connecteur
    vendor) ;
  \item le \textbf{frontend} (dashboard Next.js, authentification, couche
    données React Query) ;
  \item les \textbf{bibliothèques partagées} (domain, vault, policy-engine,
    browser-runtime, audit) ;
  \item l'\textbf{infrastructure} locale (Docker Compose, PostgreSQL,
    Redis, MinIO, Mailpit).
\end{itemize}

\section{Étude de l'existant}

\subsection{Présentation des solutions existantes}

\begin{table}[H]
\centering
\caption{Analyse comparative des approches existantes}
\label{tab:existant}
\begin{tabularx}{\textwidth}{@{}lXXC{2.4cm}@{}}
\toprule
\textbf{Solution} & \textbf{Points forts} & \textbf{Limitations}
  & \textbf{Gouvernance IA} \\
\midrule
\rowcolor{rowgray}
RPA classique (UiPath, etc.)
  & Automatisation UI mature, intégrations entreprise
  & Pas d'identité agent, pas de politiques fines, audit limité
  & Faible \\
Intégrations API directes
  & Rapides, déterministes, testables
  & Couverture partielle des portails, pas d'actions UI
  & Moyenne \\
\rowcolor{rowgray}
Gestionnaires de secrets (Vault)
  & Chiffrement robuste, rotation, audit
  & Pas de runtime navigateur, pas d'approbation workflow
  & Partielle \\
Agents IA autonomes (bruts)
  & Flexibilité maximale, intelligence adaptative
  & Risque credential, actions silencieuses, pas de preuves
  & Inexistante \\
\rowcolor{rowgray}
\textbf{AegisWeb (proposé)}
  & Identité agent, vault, politiques, approbation, preuves, audit
  & MVP local, pas encore SaaS production multi-tenant
  & Élevée \\
\bottomrule
\end{tabularx}
\end{table}

\subsection{Critique de l'existant}

L'analyse révèle plusieurs lacunes fondamentales :
\begin{itemize}[leftmargin=1.8cm, itemsep=4pt]
  \item \textbf{Absence d'identité non-humaine dédiée} pour les agents,
    avec statut, purpose et grants explicites.
  \item \textbf{Exposition des credentials} aux agents ou aux prompts LLM,
    sans chiffrement au repos ni déchiffrement scopé par exécution.
  \item \textbf{Absence de moteur de politiques} évaluant domaines,
    actions interdites, seuils financiers et niveaux de risque.
  \item \textbf{Pas de pause pour approbation humaine} avant les actions
    destructrices (downgrade, annulation).
  \item \textbf{Manque de preuves opérationnelles}~: screenshots, factures,
    reçus de confiance et chaîne d'audit vérifiable.
  \item \textbf{Runtime navigateur non contrôlé}~: accès réseau privé,
    popups hors domaine, téléchargements non tracés.
\end{itemize}

\section{Solution proposée}

\subsection{Description d'AegisWeb}

\textbf{AegisWeb} est la couche de permission, credential, approbation,
runtime navigateur et audit pour les agents IA agissant sur le web. Elle
donne aux agents une \textbf{autorité contrôlée} via~:

\begin{itemize}[leftmargin=1.8cm, itemsep=4pt]
  \item \textbf{Identité agent}~: agent non-humain avec identifier unique,
    statut (active, paused, revoked) et purpose documenté.
  \item \textbf{Coffre-fort de credentials}~: chiffrement AES-256-GCM au
    repos, octroi par agent via \texttt{CredentialAgentGrant}.
  \item \textbf{Moteur de politiques}~: allowlist de domaines, actions
    interdites, seuils de dépense, règles d'approbation.
  \item \textbf{Runtime navigateur contrôlé}~: Playwright avec allowlist,
    blocage réseau privé, masquage des champs sensibles.
  \item \textbf{Workflows queue-backed}~: trois templates SaaS procurement.
  \item \textbf{Approbation humaine}~: pause du run, notification email,
    reprise conditionnelle.
  \item \textbf{Audit trail}~: événements append-only avec chaîne de hash
    SHA-256 (\texttt{prevHash} + \texttt{eventHash}).
  \item \textbf{Receipts}~: artefacts finaux de confiance résumant
    timeline, décisions et preuves.
  \item \textbf{Dashboard opérationnel}~: salle de preuves pour agents,
    vendors, policies, runs, approvals, audit.
\end{itemize}

\begin{figure}[H]
\centering
\includegraphics[width=0.55\textwidth]{assets/aegis_agent_identity_passport.png}
\caption{Concept d'identité agent AegisWeb — passeport d'autorité contrôlée}
\label{fig:agent_passport}
\end{figure}

\subsection{Flux produit phare — Downgrade de plan SaaS}

\begin{enumerate}[leftmargin=1.8cm, itemsep=3pt]
  \item L'owner se connecte au dashboard.
  \item Il crée ou sélectionne agent, vendor, credential, policy, workflow.
  \item Il lance «~Acme Downgrade Request~».
  \item L'API crée un \texttt{WorkflowRun} et enqueue un job BullMQ.
  \item Le worker ouvre un navigateur contrôlé (Playwright).
  \item Le worker se connecte au sandbox via credential déchiffré (API
    interne vault, scopé au run).
  \item Le worker lit les données de renouvellement et de plan.
  \item Le moteur de politique classe l'action comme risquée
    $\rightarrow$ \textbf{pause pour approbation}.
  \item L'approbateur examine les preuves (screenshots) et approuve ou
    rejette.
  \item Le worker reprend uniquement après approbation.
  \item L'action est soumise $\rightarrow$ \textbf{receipt} + \textbf{audit
    trail} générés.
  \item Notification email via Mailpit (en local).
\end{enumerate}

\subsection{Objectifs du projet}

\paragraph{Objectifs fonctionnels :}
\begin{itemize}[leftmargin=1.8cm, itemsep=3pt]
  \item Développer une API REST NestJS avec 20+ modules (auth, RBAC,
    agents, vendors, policies, credentials, workflows, runs, approvals,
    receipts, audit, files, queue, internal worker).
  \item Implémenter un schéma Prisma multi-tenant (17 tables) avec seed
    de démonstration riche.
  \item Concevoir un moteur de politiques pur TypeScript sans dépendance
    BDD.
  \item Réaliser un coffre-fort AES-256-GCM avec redaction récursive des
    secrets.
  \item Développer un worker BullMQ + Playwright avec trois workflows
    SaaS.
  \item Mettre en place un dashboard Next.js 16 evidence-first.
  \item Fournir un vendor sandbox pour démos et tests déterministes.
\end{itemize}

\paragraph{Objectifs non fonctionnels :}
\begin{itemize}[leftmargin=1.8cm, itemsep=3pt]
  \item Isolation multi-tenant stricte par \texttt{organizationId}.
  \item Argon2id, JWT courte durée, refresh HttpOnly atomique.
  \item Rate limiting Redis en production.
  \item CSP nonce-based sans \texttt{unsafe-eval} sur le frontend.
  \item Tests d'intégration Vitest couvrant les 30 phases backend.
  \item Déploiement Docker Compose reproductible.
\end{itemize}

\subsection{Périmètre du projet}

Le périmètre couvre~:
\begin{enumerate}[leftmargin=1.8cm, itemsep=3pt]
  \item Le monorepo \texttt{agentpass} (pnpm workspaces).
  \item Les applications \texttt{api}, \texttt{worker}, \texttt{web},
    \texttt{vendor-sandbox}.
  \item Les bibliothèques \texttt{domain}, \texttt{database}, \texttt{vault},
    \texttt{policy-engine}, \texttt{browser-runtime}, \texttt{audit},
    \texttt{testing}.
  \item L'infrastructure \texttt{infra/docker-compose.yml}.
  \item 33 fichiers de tests Vitest + scripts smoke/E2E.
\end{enumerate}

\section{Méthodologie de gestion de projet}

\subsection{Approche incrémentale par phases}

Contrairement à une approche SCRUM classique par sprints, AegisWeb a été
développé selon une \textbf{méthodologie incrémentale en 30 phases}
(documentées dans \texttt{docs/BACKEND\_IMPLEMENTATION\_PLAN.md}), chaque
phase livrant un module testable et validé avant passage à la suivante.

\begin{table}[H]
\centering
\caption{Principes de la méthodologie par phases}
\label{tab:methodo}
\begin{tabularx}{\textwidth}{@{}lX@{}}
\toprule
\textbf{Élément}        & \textbf{Application dans le projet AegisWeb} \\
\midrule
\rowcolor{rowgray}
Phase                   & Incrément fonctionnel avec tests dédiés
                          (\texttt{phaseN.spec.ts}) \\
Definition of Done      & Tests passants, module intégré, notes de
                          complétion documentées \\
\rowcolor{rowgray}
Acceptance MVP          & \texttt{mvp-backend-acceptance.spec.ts}
                          (31 fichiers, 195 tests) \\
Traçabilité             & \texttt{docs/PHASE\_*\_COMPLETION\_NOTES.md}
                          (27 fichiers) \\
\rowcolor{rowgray}
Sécurité continue       & \texttt{security-coverage.spec.ts} +
                          \texttt{docs/SECURITY\_AUDIT.md} \\
\bottomrule
\end{tabularx}
\end{table}

\subsection{Planification générale — Diagramme de Gantt}

\begin{figure}[H]
\centering
\caption{Diagramme de Gantt — Planning prévisionnel du projet (24 semaines)}
\label{fig:gantt}
\resizebox{\textwidth}{!}{%
\begin{ganttchart}[
  hgrid,
  vgrid          = {*{1}{dotted}},
  bar/.style     = {fill=aegisblue!70, rounded corners=1pt},
  bar label font = \small,
  title label font = \small\bfseries,
  milestone label font = \small\itshape,
  milestone/.style  = {fill=aegisaccent},
  x unit         = 0.45cm,
  y unit title   = 0.62cm,
  y unit chart   = 0.68cm,
  title height   = 1,
  bar height     = 0.55,
]{1}{24}
  \gantttitle{Planning prévisionnel — AegisWeb (24 semaines)}{24} \\
  \gantttitlelist{1,...,24}{1} \\
  \ganttbar{Étude \& Analyse besoins}{1}{2} \\
  \ganttbar{Phases 0--3 : Fondation API}{3}{5} \\
  \ganttbar{Phases 4--8 : Auth, RBAC, Org, Users}{6}{8} \\
  \ganttbar{Phases 9--13 : Agents, Vendors, Vault}{9}{11} \\
  \ganttbar{Phases 14--18 : Workflows, Queue}{12}{14} \\
  \ganttbar{Phases 19--23 : Worker, Browser, Sandbox}{15}{17} \\
  \ganttbar{Phases 24--26 : Workflows SaaS}{18}{19} \\
  \ganttbar{Phases 27--29 : Receipts, Notif, OpenAPI}{20}{21} \\
  \ganttbar{Dashboard Next.js (frontend)}{10}{22} \\
  \ganttbar{Tests acceptance \& sécurité}{21}{23} \\
  \ganttbar{Rédaction rapport PFE}{22}{24} \\
  \ganttmilestone{Livraison finale}{24}
\end{ganttchart}%
}
\end{figure}

\conclusionchap{Ce chapitre nous a permis de situer AegisWeb dans le
contexte de la gouvernance des agents IA web, d'identifier les
insuffisances des solutions existantes et de définir les contours de
notre plateforme. La méthodologie incrémentale en 30 phases et le planning
sur 24 semaines constituent le cadre organisationnel de notre démarche.
Le chapitre suivant sera consacré à l'analyse détaillée des besoins.}
